\ifdefined\gkver\else\def\gkver{student}\fi
\documentclass[\gkver]{gaokaozhenti}
\begin{document}

\chapter{2017年11月浙江}

\begin{enumerate}
\item
%% number: 1
%% paperName: 2017年11月浙江选考第 1 题 3 分
%% typeId: 1
%% type: 单选题
%% chapter: 直线运动
%% point: 路程与位移
%% method: 无
%% score: 3
%% degree: 950
%% duplicateId: 0
%% body:
下列物理量中属于矢量的是\xzanswer{D}
\fourchoices[answer=D]
{速率}
{电势}
{电流}
{位移}
%% answer: D
%% memo:
\memoanswer{
	矢量既有大小、又有方向，且运算遵从平行四边形定则，而标量只有大小、没有方向，运算遵从代数运算法则。瞬时速度的大小称之为速率，是标量，故 A 错误；电势是只有大小没有方向的标量，故 B 错误；电流虽有方向，但其方向只表示正、负，且运算遵从代数运算法则，是标量，故 C 错误；位移是既有大小又有方向的矢量，故 D 正确。
}

\item
%% number: 2
%% paperName: 2017年11月浙江选考第 2 题 3 分
%% typeId: 1
%% type: 单选题
%% chapter: 运动和力
%% point: 力学单位制
%% method: 无
%% score: 3
%% degree: 950
%% duplicateId: 0
%% body:
在国际单位制中，属于基本量及基本单位的是\xzanswer{A}
\fourchoices[answer=A]
{质量，千克}
{能量，焦耳}
{电阻，欧姆}
{电量，库仑}
%% answer: A
%% memo:
\memoanswer{
	国际单位制中基本物理量有：质量、长度、时间、电流、发光强度、热力学温度、物质的量，其基本单位分别为千克 $\Ukg$、米 $\Um$、秒 $\Us$、安培 $\UA$、坎德拉 $\Ucd$、开尔文 $\UK$、摩尔 $\Umol$，故 A 正确。
}

\item
%% number: 3
%% paperName: 2017年11月浙江选考第 3 题 3 分
%% typeId: 1
%% type: 单选题
%% chapter: 其他
%% point: 物理学史
%% method: 无
%% score: 3
%% degree: 950
%% duplicateId: 0
%% body:
获得 2017 年诺贝尔物理学奖的成果是\xzanswer{D}
\fourchoices[answer=D]
{牛顿发现了万有引力定律}
{卡文迪许测定了引力常量}
{爱因斯坦预言了引力波}
{雷纳$\cdot$韦斯等探测到了引力波}
%% answer: D
%% memo:
\memoanswer{
	2017 年最热门的物理学发现是引力波，故 D 正确。
}

\item
%% number: 4
%% paperName: 2017年11月浙江选考第 4 题 3 分
%% typeId: 1
%% type: 单选题
%% chapter: 直线运动
%% point: 路程与位移
%% method: 无
%% score: 3
%% degree: 750
%% duplicateId: 0
%% body:
如图所示，两位同学从滑道最高端的同一位置先后滑下，到达底端的同一位置，对于整个下滑过程，两同学的\xzanswer{A}

\onepicture[w=3.5cm]{figs/04.jpg}

\fourchoices[answer=A]
{位移一定相同}
{时间一定相同}
{末速度一定相同}
{平均速度一定相同}
%% answer: A
%% memo:
\memoanswer{
	两位同学初、末位置相同，即两位同学位置的变化是一样的，故两位同学的位移相同，故 A 正确；两位同学滑下来的时间不一定相同，在轨道上所受的摩擦力也不一定相同，故两位同学的末速度、平均速度不一定相同，故 BCD 错误。
}

\item
%% number: 5
%% paperName: 2017年11月浙江选考第 5 题 3 分
%% typeId: 1
%% type: 单选题
%% chapter: 力物体的平衡
%% point: 多力平衡
%% method: 整体与隔离
%% score: 3
%% degree: 650
%% duplicateId: 0
%% body:
叠放在水平地面上的四个完全相同的排球如图所示，质量均为 $m$，相互接触，球与地面间的动摩擦因数均为 $\mu$，则\xzanswer{C}

\onepicture[w=4.5cm]{figs/05.png}

\fourchoices[answer=C]
{上方球与下方 3 个球间均没有弹力}
{下方三个球与水平地面间均没有摩擦力}
{水平地面对下方三个球的支持力均为 $\frac{4}{3}mg$}
{水平地面对下方三个球的摩擦力均为 $\frac{4}{3}\mu mg$}
%% answer: C
%% memo:
\memoanswer{
	将四个球看作一个整体，地面的支持力与球重力之和平衡，设其中一个球受的支持力大小为 $F_N$，因此 $3F_N=4mg$，即 $F_N=\frac{4}{3}mg$，故 C 正确；经受力分析，最上面的球对下面三个球肯定有弹力作用，故 A 错误；对地面上的其中一个球受力分析，如图所示，可知球受到的摩擦力是地面的静摩擦力，因此不能通过 $F_f=\mu F_N$ 求解，故 BD 错误。

	\onepicture[w=5.0cm]{figs/05_an.jpg}
}

\item
%% number: 6
%% paperName: 2017年11月浙江选考第 6 题 3 分
%% typeId: 1
%% type: 单选题
%% chapter: 电场
%% point: 电场线
%% method: 无
%% score: 3
%% degree: 750
%% duplicateId: 0
%% body:
电场线的形状可以用实验来模拟，把头发屑悬浮在蓖麻油里，加上电场，头发屑就按照电场的方向排列起来，如图所示，关于此实验，下列说法正确的是\xzanswer{C}

\twopicture[labelA=201711月浙江06a, labelB=201711月浙江06b, w=5.2cm]
{figs/06a.png}
{figs/06b.png}

\fourchoices[answer=C]
{\ref{201711月浙江06a} 图是模拟两等量同种电荷的电场线}
{\ref{201711月浙江06b} 图一定是模拟两等量正电荷的电场线}
{\ref{201711月浙江06a} 图中 A、B 应接高压起电装置的两极}
{\ref{201711月浙江06b} 图中 A、B 应接高压起电装置的两极}
%% answer: C
%% memo:
\memoanswer{
	图 \ref{201711月浙江06a} 为模拟异种电荷形成的电场线，但无法分辨是否为等量异种电荷形成的电场线分布；图 \ref{201711月浙江06b} 为同种电荷，但也无法分辨是否为等量正电荷形成的电场线分布；图 \ref{201711月浙江06a} 中 A、B 相吸引，图 \ref{201711月浙江06b} 中 A、B 相排斥，故图 \ref{201711月浙江06a} 中 A、B 接的是高压起电装置的两极，故 C 正确。
}

\item
%% number: 7
%% paperName: 2017年11月浙江选考第 7 题 3 分
%% typeId: 1
%% type: 单选题
%% chapter: 曲线运动
%% point: 人造地球卫星
%% method: 无
%% score: 3
%% degree: 750
%% duplicateId: 0
%% body:
如图所示是小明同学画的人造卫星轨道的示意图，则卫星\xzanswer{D}

\onepicture[w=4.5cm]{figs/07.png}

\fourchoices[answer=D]
{在 a 轨道运动的周期为 $24\Uh$}
{在 b 轨道运动的速度始终不变}
{在 c 轨道运行的速度大小始终不变}
{在 c 轨道运行时受到的地球引力大小是变化的}
%% answer: D
%% memo:
\memoanswer{
	同步卫星必须在赤道正上空 $36000\Ukm$ 处，所以 a 轨道不可能是同步轨道，故 A 错误；b 轨道上的卫星的速度方向不断变化，即速度在变，故 B 错误；地球在 c 轨道的其中一个焦点上，因此在近地点卫星速度快，远地点速度慢，故 C 错误；在 c 轨道上，卫星离地球的距离在变化，由 $F=G\frac{Mm}{r^{2}}$ 可知卫星受地球的引力大小不断变化，故 D 正确。
}

\item
%% number: 8
%% paperName: 2017年11月浙江选考第 8 题 3 分
%% typeId: 1
%% type: 单选题
%% chapter: 磁场
%% point: 带电粒子在复合场中的运动
%% method: 无
%% score: 3
%% degree: 750
%% duplicateId: 0
%% body:
如图所示，在两水平金属板构成的器件中，存在着匀强电场与匀强磁场，电场强度 $E$ 和磁感应强度 $B$ 相互垂直，以某一水平速度进入的不计重力的带电粒子恰好能沿直线运动，下列说法正确的是\xzanswer{C}

\onepicture[w=5.0cm]{figs/08.png}

\fourchoices[answer=C]
{粒子一定带负电}
{粒子的速度大小 $v=\frac{B}{E}$}
{若粒子速度大小改变，粒子将做曲线运动}
{若粒子速度大小改变，电场对粒子的作用力会发生改变}
%% answer: C
%% memo:
\memoanswer{
	粒子做直线运动，说明竖直方向受力平衡，即 $qvB=qE$，得 $v=\frac{E}{B}$，故 B 错误；假设粒子带正电荷，则洛伦兹力竖直向上，电场力竖直向下，若粒子带负电，则二力方向相反，只要满足 $v=\frac{E}{B}$ 依然可以满足题意，故 A 错误；粒子速度变大或变小，都会导致洛伦兹力变化，因此粒子会做曲线运动，故 C 正确；不管粒子速度怎么变，只要粒子所带电量不变，在匀强电场中，粒子受到的电场力不变，故 D 错误。
}

\item
%% number: 9
%% paperName: 2017年11月浙江选考第 9 题 3 分
%% typeId: 1
%% type: 单选题
%% chapter: 直线运动
%% point: 运动图象
%% method: 无
%% score: 3
%% degree: 650
%% duplicateId: 0
%% body:
杂技运动在训练时的照片如图所示，有一小球自由落下，落到水平桌面后反弹，如此数次落下和反弹。若规定竖直向下为正方向，碰撞时间不计，空气阻力不计，则下列 $v-t$ 图像中正确的是\xzanswer{B}

\onepicture[w=3.0cm]{figs/09a.jpg}

\fourchoices[answer=B, ispicture=true, h=1.8cm]
{figs/09b.png}
{figs/09c.png}
{figs/09d.png}
{figs/09e.png}
%% answer: B
%% memo:
\memoanswer{
	小球先向下做自由落体运动，碰到桌面反弹起来，所以速度先是正值，后变为负值，且小球在弹起来时速度大小变化很小，仅改变方向，故 B 正确。
}

\item
%% number: 10
%% paperName: 2017年11月浙江选考第 10 题 3 分
%% typeId: 1
%% type: 单选题
%% chapter: 功和能
%% point: 功率
%% method: 无
%% score: 3
%% degree: 550
%% duplicateId: 0
%% body:
如图所示，质量为 $60\Ukg$ 的某运动员在做俯卧撑运动，运动过程中可将她的身体视为一根直棒，已知重心在 c 点，其垂线与脚、两手连线中点间的距离 oa、ob 分别为 $0.9\Um$ 和 $0.6\Um$，若她在 $1\Umin$ 内做了 30 个俯卧撑，每次肩部上升的距离均为 $0.4\Um$，则克服重力做功和相应的功率为\xzanswer{B}

\onepicture[w=4.5cm]{figs/10.png}

\fourchoices[answer=B]
{$430\UJ$，$7\UW$}
{$4300\UJ$，$70\UW$}
{$720\UJ$，$12\UW$}
{$7200\UJ$，$120\UW$}
%% answer: B
%% memo:
\memoanswer{
	由相似三角形知，每次运动员做俯卧撑中，重心变化的高度 $\frac{h}{0.4}=\frac{0.9}{0.9+0.6}$，即 $h=0.24\Um$。一次俯卧撑中，克服重力做的功为 $W=mgh=60\times10\times0.24\UJ=144\UJ$，所以一分钟内克服重力做功总共 $W_{\text{总}}=NW=4320\UJ$，功率 $P=\frac{W_{\text{总}}}{t}=72\UW$，故 B 正确。
}

\item
%% number: 11
%% paperName: 2017年11月浙江选考第 11 题 3 分
%% typeId: 1
%% type: 单选题
%% chapter: 曲线运动
%% point: 圆周运动的应用
%% method: 无
%% score: 3
%% degree: 750
%% duplicateId: 0
%% body:
如图所示，照片中的汽车在水平路面上做匀速圆周运动，已知图中双向四车道的总宽度约为 $15\Um$，假设汽车受到的最大静摩擦力等于车重的 0.7 倍，则运动的汽车\xzanswer{C}

\onepicture[w=5.0cm]{figs/11.jpg}

\fourchoices[answer=C]
{所受的合力可能为零}
{只受重力和地面的支持力作用}
{最大速度不能超过 $25\Ums$}
{所需的向心力由重力和支持力的合力提供}
%% answer: C
%% memo:
\memoanswer{
	汽车在水平面上做匀速圆周运动，合外力时刻指向圆心，即汽车所受的合外力不为 0；拐弯时由静摩擦力提供向心力，因此排除选项 ABD，故 C 正确。
}

\item
%% number: 12
%% paperName: 2017年11月浙江选考第 12 题 3 分
%% typeId: 1
%% type: 单选题
%% chapter: 电路
%% point: 电功电功率
%% method: 无
%% score: 3
%% degree: 750
%% duplicateId: 0
%% body:
小明同学家里部分电器的消耗功率及每天工作时间如表所示，则这些电器一天消耗的电能约为\xzanswer{D}

\begin{center}
\begin{tabular}{|c|c|c|}
\hline
电器 & 消耗功率/W & 工作时间/h \\
\hline
电茶壶 & 2000 & 1 \\
\hline
空调 & 1200 & 3 \\
\hline
电视机 & 100 & 2 \\
\hline
节能灯 & 16 & 4 \\
\hline
路由器 & 9 & 24 \\
\hline
\end{tabular}
\end{center}

\fourchoices[answer=D]
{$6.1\times10^{3}\UW$}
{$6.1\times10^{3}\UJ$}
{$2.2\times10^{4}\UW$}
{$2.2\times10^{7}\UJ$}
%% answer: D
%% memo:
\memoanswer{
	用电器消耗的电能 $W=Pt$，则由表中数据，每天消耗的电能为
	$W=2\UkW\times1\Uh+1.2\UkW\times3\Uh+0.1\UkW\times2\Uh+0.016\UkW\times4\Uh+0.009\UkW\times24\Uh=6.08\UkWh\approx2.2\times10^{7}\UJ$，故 D 正确。
}

\item
%% number: 13
%% paperName: 2017年11月浙江选考第 13 题 3 分
%% typeId: 1
%% type: 单选题
%% chapter: 功和能
%% point: 功率
%% method: 无
%% score: 3
%% degree: 550
%% duplicateId: 0
%% body:
如图所示是具有登高平台的消防车，具有一定质量的伸缩臂能够在 $5\Umin$ 内使承载 4 人的登高平台（人连同平台的总质量为 $400\Ukg$）上升 $60\Um$ 到达灭火位置，此后，在登高平台上的消防员用水炮灭火，已知水炮的出水量为 $3\Umc/\Umin[a]$，水离开炮口时的速率为 $20\Ums$，则用于\xzanswer{B}

\onepicture[w=5.5cm]{figs/13.jpg}

\fourchoices[answer=B]
{水炮工作的发动机输出功率约为 $1\times10^{4}\UW$}
{水炮工作的发动机输出功率约为 $4\times10^{4}\UW$}
{水炮工作的发动机输出功率约为 $2.4\times10^{6}\UW$}
{伸缩臂抬升登高平台的发动机输出功率约为 $800\UW$}
%% answer: B
%% memo:
\memoanswer{
	伸缩臂抬升登高平台的发动机输出功率为设备升高克服重力的功率，即 $P=\frac{W}{t}$，由于伸缩臂也要升高，但其质量未知，故不能计算抬升登高平台的发动机的功率，故 D 错误；在一秒钟内，喷出去水的质量为 $m=\rho V=10^{3}\times\frac{1}{20}\Ukg=50\Ukg$，喷出去水的重力势能为 $W_G=mgh=50\times10\times60\UJ=3\times10^{4}\UJ$，水的动能为 $\frac{1}{2}mv^{2}=1\times10^{4}\UJ$，所以一秒钟内水增加的能量为 $4\times10^{4}\UJ$，即水炮工作的发动机功率约为 $4\times10^{4}\UW$，故 B 正确，AC 错误。
}

\item
%% number: 14
%% paperName: 2017年11月浙江选考第 14 题 2 分
%% typeId: 2
%% type: 多选题
%% chapter: 原子和原子核
%% point: 核裂变
%% method: 无
%% score: 2
%% degree: 650
%% duplicateId: 0
%% body:
下列说法正确的是\xzanswer{BD}
\fourchoices[answer=BD]
{核聚变反应方程可能为 \ce{^2_1H + ^3_1H -> ^4_2He + 2^1_0n}}
{铀核裂变的核反应方程可能为 \ce{^235_92U + ^1_0n -> ^140_54Xe + ^94_38Sr + 2^1_0n}}
{发生 $\beta$ 衰变时原子核发出电子，说明电子是原子核的组成部分}
{中子和质子结合成氘核，若该过程质量亏损为 $\Delta m$，则氘核的结合能为 $\Delta mc^{2}$}
%% answer: BD
%% memo:
\memoanswer{
	选项 A 的核反应方程应为 \ce{^2_1H + ^3_1H -> ^4_2He + ^1_0n}，故 A 错误；由核反应前后质量数、质子数守恒可知，选项 B 正确；$\beta$ 衰变是原子核内部中子衰变为质子并释放出高速电子，电子并不是原子核的组成部分，故 C 错误；由质能方程可知，选项 D 正确。
}

\item
%% number: 15
%% paperName: 2017年11月浙江选考第 15 题 2 分
%% typeId: 2
%% type: 多选题
%% chapter: 光学
%% point: 光的干涉
%% method: 无
%% score: 2
%% degree: 450
%% duplicateId: 0
%% body:
a、b 是两种单色光，其光子能量分别为 $\varepsilon_a$ 和 $\varepsilon_b$，且 $\frac{\varepsilon_a}{\varepsilon_b}=k$，\xzanswer{BC}
\fourchoices[answer=BC]
{则 a、b 的光子动量之比 $\frac{p_a}{p_b}=1:k$}
{若 a、b 入射到同一双缝干涉装置上，则相邻亮条纹的间距之比 $\frac{\Delta x_a}{\Delta x_b}=1:k$}
{若 a、b 都能使某种金属发生光电效应，则光电子最大初动能之差 $E_{ka}-E_{kb}=\varepsilon_b(k-1)$}
{若 a、b 是由处于同一激发态的原子跃迁到 a 态和 b 态时产生的，则 a、b 两态能级之差 $E_a-E_b=\varepsilon_b(k-1)$}
%% answer: BC
%% memo:
\memoanswer{
	由德布罗意方程 $\lambda=\frac{h}{p}$ 以及光子的能量 $\varepsilon=h\frac{c}{\lambda}$，联立方程 $p=\frac{\varepsilon}{c}$，所以 $\frac{p_a}{p_b}=\frac{\varepsilon_a}{\varepsilon_b}=k$，故 A 错误；$\Delta x=\frac{L\lambda}{d}=\frac{Lhc}{d\varepsilon}$，故 B 正确；光电子最大初动能 $E_k=\varepsilon-W_0$，因此 $E_{ka}-E_{kb}=(k-1)\varepsilon_b$，故 C 正确；能级跃迁如图所示，$E_a-E_b=\varepsilon_b-\varepsilon_a=\varepsilon_b(1-k)$，故 D 错误。

	\onepicture[w=6.5cm]{figs/15_an.jpg}
}

\item
%% number: 16
%% paperName: 2017年11月浙江选考第 16 题 2 分
%% typeId: 2
%% type: 多选题
%% chapter: 机械波
%% point: 干涉衍射
%% method: 无
%% score: 2
%% degree: 450
%% duplicateId: 0
%% body:
有两列频率相同、振动方向相同、振幅均为 $A$、传播方向互相垂直的平面波相遇发生干涉。如图所示，图中实线表示波峰，虚线表示波谷，a 为波谷与波谷相遇点，b、c 为波峰与波谷相遇点，d 为波峰与波峰相遇点，e、g 是 a、d 连线上的两点，其中 e 为连线的中点，则\xzanswer{ABC}

\onepicture[w=5.0cm]{figs/16.png}

\fourchoices[answer=ABC]
{a、d 处的质点振动加强，b、c 处的质点振动减弱}
{从图示时刻经过半个周期，e 处的质点通过的路程为 $4A$}
{从图示时刻经过半个周期，g 处的质点加速向平衡位置运动}
{从图示时刻经过四分之一周期，d 处的质点振幅恰好为零}
%% answer: ABC
%% memo:
\memoanswer{
	由题意可知两列波满足发生干涉的条件，a、d 两点分别为波谷与波谷叠加和波峰与波峰叠加，属于振动加强点，而 b、c 两点为波峰与波谷叠加，属于振动减弱点，故 A 正确；此刻传播的机械波，对应的 e、g 两点分别在两列波中的位置以及振动情况如图，经过半个周期后，e 点的路程为全振动路程的一半，故 B 正确；经半个周期后，由图中可以看出 g 点的合运动为向下加速靠近平衡位置，故 C 正确；在机械波中质点振动的振幅为偏离平衡位置的最大距离，是一定值，不会发生变化，所以 d 点的振幅为 $2A$，故 D 错误。

	\onepicture[w=6.5cm]{figs/16_an.jpg}
}

\item
%% number: 17
%% paperName: 2017年11月浙江选考第 17 题 5 分
%% typeId: 4
%% type: 实验题
%% chapter: 运动和力
%% point: 验证牛顿第二定律
%% method: 无
%% score: 5
%% degree: 650
%% duplicateId: 0
%% body:
在做“探究加速度与力、质量的关系”实验中，
\begin{enumerate}
	\item
	下列仪器需要用到的是 \tkanswer{AD}（多选）；

	\fourchoices[answer=AD, ispicture=true, h=3.0cm]
	{figs/17a.png}
	{figs/17b.png}
	{figs/17c.png}
	{figs/17d.png}

	\item
	下列说法正确的是 \tkanswer{BD}（多选）；

	（A）先释放纸带再接通电源

	（B）拉小车的细线应尽可能与长木板平行

	（C）纸带与小车相连端的点迹较疏

	（D）轻推小车，拖着纸带的小车能够匀速下滑说明摩擦力已被平衡

	\item
	如图所示是实验时打出的一条纸带，A、B、C、D……为每隔 4 个点取的计数点，据此纸带可知小车在 D 点速度大小为 \tkanswer{0.21}$\Ums$（小数点后保留两位）。

	\onepicture[h=2.0cm, align=right]{figs/17e.png}
\end{enumerate}
%% answer: 
\jdanswer{
\begin{enumerate}
	\item
	AD

	\item
	BD

	\item
	0.21

\end{enumerate}
}
%% memo:
\memoanswer{
	（1）本实验通过分析打点纸带上点的分布求物体的加速度，需要通过改变小车质量，利用控制变量法研究加速度、质量、合外力之间的关系，所以选择仪器 AD。

	（2）实验时应先打开打点计时器电源，后释放小车，为了保证绳子拉力不用分解，绳子要与木板平行，故 A 错误 B 正确；小车释放前应尽量靠近打点计时器，以保证点足够多，便于分析小车加速度，且初始的小车的速度较小，所以纸带与小车相连端点迹密，故 C 错误；本实验需要平衡摩擦力，在小车质量远大于钩码质量时，能将钩码重力近似当作小车合外力，故 D 正确。

	（3）$t=5T=0.1\Us$，则 $v_D=\frac{x_{CE}}{2t}=0.205\Ums$，保留小数点后两位，所以答案为 $0.21\Ums$。
}

\item
%% number: 18
%% paperName: 2017年11月浙江选考第 18 题 5 分
%% typeId: 4
%% type: 实验题
%% chapter: 电路
%% point: 测电动势与内阻
%% method: 无
%% score: 5
%% degree: 550
%% duplicateId: 0
%% body:
小明同学在测定一节干电池的电动势和内阻的实验时，为防止电流过大而损坏器材，电路中加了一个保护电阻 $R_0$，根据如图 \ref{201711月浙江18a} 所示电路图进行实验时，
\begin{enumerate}
	\item
	电流表量程应选择 \tkanswer{$0.6\UA$}（填“$0.6\UA$”或“$3\UA$”），保护电阻应选用 \tkanswer{B}（填“A”或“B”）；

	（A）定值电阻（阻值 $10.0\UO$，额定功率 $10\UW$）

	（B）定值电阻（阻值 $2.0\UO$，额定功率 $5\UW$）

	\item
	在一次测量中电压表的指针位置如图 \ref{201711月浙江18b} 所示，则此时电压为 \tkanswer{1.21}$\UV$；

	\item
	根据实验测得的 5 组数据画出的 $U-I$ 图线如图 \ref{201711月浙江18c} 所示，则干电池的电动势 $E=$ \tkanswer{1.45}$\UV$，内阻 $r=$ \tkanswer{0.50}$\UO$（小数点后保留两位）。
\end{enumerate}

\threepicture[numA=1, numB=2, numC=3, numstyle=arabic, labelA=201711月浙江18a, labelB=201711月浙江18b, labelC=201711月浙江18c, align=right, wA=4.0cm, wB=4.2cm, wC=4.2cm]
{figs/18a.png}
{figs/18b.jpg}
{figs/18c.png}
%% answer: 
\jdanswer{
\begin{enumerate}
	\item
	0.6 A；B

	\item
	1.21

	\item
	1.45；0.50

\end{enumerate}
}
%% memo:
\memoanswer{
	（1）通过题图丙的 $U-I$ 图象可知，电流表选择 $0.6\UA$ 即可。而电路中的电流最大时 $I=\frac{E}{R_0}=0.6\UA$，可求得保护电阻 $R_0=2.5\UO$，若选择 $10\UO$ 的定值电阻，则短路电流不足以满足偏转量程的 $\frac{1}{3}$，因此选择 B 比较合理。

	（2）根据读数方法可知电压表读数约 $1.20\UV$。

	（3）实验中，保护电阻可以看作电源内阻一部分，所以通过图象可知 $E=1.45\UV$，$r+R_0=\frac{1.45-0.70}{0.30}\UO=2.5\UO$，所以内阻 $r=0.5\UO$。
}

\item
%% number: 19
%% paperName: 2017年11月浙江选考第 19 题 9 分
%% typeId: 6
%% type: 计算题
%% chapter: 曲线运动
%% point: 平抛运动
%% method: 无
%% score: 9
%% degree: 550
%% duplicateId: 0
%% body:
如图所示，AMB 是一条长 $L=10\Um$ 的绝缘水平轨道，固定在离水平地面高 $h=1.25\Um$ 处，A、B 为端点，M 为中点，轨道 MB 处在方向竖直向上、大小 $E=5\times10^{3}\UNC$ 的匀强电场中，一质量 $m=0.1\Ukg$，电荷量 $q=+1.3\times10^{-4}\UC$ 的可视为质点的滑块以初速度 $v_0=6\Ums$ 在轨道上自 A 点开始向右运动，经 M 点进入电场，从 B 点离开电场，已知滑块与轨道间动摩擦因数 $\mu=0.2$，求滑块
\begin{enumerate}
	\item
	到达 M 点时的速度大小；
	\item
	从 M 点运动到 B 点所用的时间；
	\item
	落地点距 B 点的水平距离。
\end{enumerate}

\onepicture[w=7.0cm, align=right]{figs/19.png}
%% answer: 
\jdanswer{
\begin{enumerate}
	\item
	$4\Ums$

	\item
	$\frac{10}{7}\Us$

	\item
	$1.5\Um$

\end{enumerate}
}
%% memo:
\memoanswer{
	（1）滑块在 AM 段的加速度为 $a_1$，由牛顿第二定律得

	$a_1=\frac{F_f}{m}=\frac{\mu mg}{m}=2\Umsq$，

	由运动学公式得

	$v_M^{2}-v_A^{2}=2(-a_1)\frac{L}{2}$，解得 $v_M=4\Ums$。

	（2）滑块在 MB 段的加速度为 $a_2$，由牛顿第二定律得

	$a_2=\frac{\mu(mg-qE)}{m}=0.7\Umsq$，

	由运动学公式得

	$\frac{L}{2}=v_Mt_1-\frac{1}{2}a_2t_1^{2}$，解得 $t_1=\frac{10}{7}\Us$。

	（3）滑块下落时间为 $t_2=\sqrt{\frac{2h}{g}}=0.5\Us$，

	由运动学公式得

	$v_B=v_M-a_2t_1=3\Ums$，

	$x=v_Bt_2=1.5\Um$。
}

\item
%% number: 20
%% paperName: 2017年11月浙江选考第 20 题 12 分
%% typeId: 6
%% type: 计算题
%% chapter: 功和能
%% point: 动能定理
%% method: 无
%% score: 12
%% degree: 450
%% duplicateId: 0
%% body:
如图 \ref{201711月浙江20a} 所示是游乐园的过山车，其局部可简化为如图 \ref{201711月浙江20b} 所示的示意图，倾角 $\theta=37^\circ$ 的两平行倾斜轨道 BC、DE 的下端与水平半圆形轨道 CD 顺滑连接，倾斜轨道 BC 的 B 端高度 $h=24\Um$，倾斜轨道 DE 与圆弧 EF 相切于 E 点，圆弧 EF 的圆心 $O_1$、水平半圆轨道 CD 的圆心 $O_2$ 与 A 点在同一水平面上，$DO_1$ 的距离 $L=20\Um$，质量 $m=1000\Ukg$ 的过山车（包括乘客）从 B 点自静止滑下，经过水平半圆轨道后，滑上另一倾斜轨道，到达圆弧顶端 F 时乘客对座椅的压力为自身重力的 0.25 倍。已知过山车在 BCDE 段运动时所受的摩擦力与轨道对过山车的支持力成正比，比例系数 $\mu=\frac{1}{32}$，EF 段摩擦力不计，整个运动过程空气阻力不计。（$\sin37^\circ=0.6$，$\cos37^\circ=0.8$）
\begin{enumerate}
	\item
	求过山车过 F 点时的速度大小；
	\item
	求从 B 到 F 整个运动过程中摩擦力对过山车做的功；
	\item
	如图过 D 点时发现圆轨道 EF 段有故障，为保证乘客的安全，立即触发制动装置，使过山车不能到达 EF 段并保证不再下滑，则过山车受到的摩擦力至少多大？
\end{enumerate}

\twopicture[numA=1, numB=2, numstyle=arabic, labelA=201711月浙江20a, labelB=201711月浙江20b, align=right, wA=5.5cm, wB=7.5cm]
{figs/20a.jpg}
{figs/20b.jpg}
%% answer: 
\jdanswer{
\begin{enumerate}
	\item
	$3\sqrt{10}\Ums$

	\item
	$W=-7.5\times10^{4}\UJ$

	\item
	$6\times10^{3}\UN$

\end{enumerate}
}
%% memo:
\memoanswer{
	（1）过山车在 F 点，由牛顿第二定律得

	$mg-\frac{1}{4}m_{\text{人}}g=m\frac{v_F^{2}}{r}$①，

	又 $r=L\sin\theta=12\Um$②，

	联立①②解得 $v_F=\sqrt{\frac{3}{4}gr}=3\sqrt{10}\Ums$③。

	（2）设整个过程摩擦阻力做功为 $W$，对 B 到 F 的过程，由动能定理得

	$mg(h-r)+W=\frac{1}{2}mv_F^{2}-0$④，

	解得 $W=-7.5\times10^{4}\UJ$⑤。

	（3）触发制动后能恰好到达 E 点对应的摩擦力为 $F_{f1}$，由动能定理得

	$-F_{f1}L\cos\theta-mgr\cos\theta=0-\frac{1}{2}mv_D^{2}$⑥，

	未触发制动时，对 D 点到 F 点的过程，由动能定理得

	$-\mu mg\cos\theta\cdot L\cos\theta-mgr=\frac{1}{2}mv_F^{2}-\frac{1}{2}mv_D^{2}$⑦，

	由⑥⑦两式得 $F_{f1}=\frac{73}{16}\times10^{3}\UN=4562.5\UN$⑧，

	要使过山车停在倾斜轨道上的摩擦力为 $F_{f2}$，由受力分析得

	$F_{f2}=mg\sin\theta=6\times10^{3}\UN$⑨，

	综合考虑⑧⑨两式，得 $F_{fm}=6\times10^{3}\UN$。
}

\item
%% number: 21
%% paperName: 2017年11月浙江选考第 21 题 4 分
%% typeId: 4
%% type: 实验题
%% chapter: 光学
%% point: 测定玻璃折射率
%% method: 无
%% score: 4
%% degree: 650
%% duplicateId: 0
%% body:
在“测定玻璃的折射率”实验时
\begin{enumerate}
	\item
	下列说法正确的是 \tkanswer{C}；

	（A）入射角越大，误差越小

	（B）在白纸上放好玻璃砖后，用铅笔贴着光学面画出界面

	（C）实验时既可用量角器，也可用圆规和直尺等工具进行测量

	（D）判断像与针是否在同一直线时，应该观察大头针的头部

	\item
	小时同学在插针时玻璃砖的位置如图 \ref{201711月浙江21a} 所示。根据插针与纸上已画的界面确定入射点与出射点，依据上述操作所测得的折射率 \tkanswer{偏小}（填“偏大”、“偏小”或“不变”）；

	\item
	小明同学经正确操作后，在纸上留下四枚大头针的位置 $P_1$、$P_2$、$P_3$、$P_4$，AB 和 CD 是描出的玻璃砖的两个边，如图 \ref{201711月浙江21b} 所示，请在答题纸上画出光路图。
\end{enumerate}

\twopicture[numA=1, numB=2, numstyle=arabic, labelA=201711月浙江21a, labelB=201711月浙江21b, align=right, wA=4.5cm, wB=3.5cm]
{figs/21a.png}
{figs/21b.png}
%% answer: 
\jdanswer{
\begin{enumerate}
	\item
	C

	\item
	偏小

	\item
	光路图如解析图所示

\end{enumerate}
}
%% memo:
\memoanswer{
	（1）本实验在用插针法时，入射角度不宜过大，否则很难确定折射角，故 A 错误；在实验前，首先应该确定玻璃砖上表面，用铅笔在纸面上画一直线，然后确定下表面，故 B 错误；本实验通过 $n=\frac{\sin i}{\sin r}$ 测量折射率，也可以通过画出两个半径一样的圆，通过 $\sin\theta=\frac{x_1}{r}$ 将角度之比转化为角度对边之比，故 C 正确；在观察像与针是否在一条直线时，应该观察大头针底部，故 D 错误。

	（2）玻璃砖上表面水平，下表面向上倾斜，作出光路图如图，由图可知利用 $n=\frac{\sin i}{\sin r}$ 计算时，测量角度偏大，导致测量折射率偏小。

	\onepicture[w=4.5cm]{figs/21_an1.jpg}

	（3）光路图如图所示。

	\onepicture[w=4.0cm]{figs/21_an2.png}
}

\item
%% number: 22
%% paperName: 2017年11月浙江选考第 22 题 10 分
%% typeId: 6
%% type: 计算题
%% chapter: 电磁感应
%% point: 电磁感应能量分析
%% method: 无
%% score: 10
%% degree: 350
%% duplicateId: 0
%% body:
如图所示，匝数 $N=100$、截面积 $S=1.0\times10^{-2}\Umq$、电阻 $r=0.15\UO$ 的线圈内有方向垂直于线圈平面向上的随时间均匀增加的匀强磁场 $B_1$，其变化率 $k=0.80\UTs$。线圈通过开关 S 连接两根相互平行、间距 $d=0.20\Um$ 的竖直导轨，下端连接阻值 $R=0.50\UO$ 的电阻。一根阻值也为 $0.50\UO$、质量 $m=1.0\times10^{-2}\Ukg$ 的导体棒 ab 搁置在等高的挡条上。在竖直导轨间的区域仅有垂直纸面的不随时间变化的匀强磁场 $B_2$。接通开关 S 后，棒对挡条的压力恰好为零。假设棒始终与导轨垂直，且与导轨接触良好，不计摩擦阻力和导轨电阻。
\begin{enumerate}
	\item
	求磁感应强度 $B_2$ 的大小，并指出磁场方向；
	\item
	断开开关 S 后撤去挡条，棒开始下滑，经 $t=0.25\Us$ 后下降了 $h=0.29\Um$，求此过程棒上产生的热量。
\end{enumerate}

\onepicture[h=4.5cm, align=right]{figs/22.png}
%% answer: 
\jdanswer{
\begin{enumerate}
	\item
	$0.50\UT$，垂直纸面向外

	\item
	$2.3\times10^{-3}\UJ$

\end{enumerate}
}
%% memo:
\memoanswer{
	（1）由法拉第电磁感应定律得，线圈产生的感应电动势为

	$E=N\frac{\Delta\Phi}{\Delta t}=NS\frac{\Delta B_1}{\Delta t}=NSk$，

	流过导体棒的电流为

	$I_{ab}=\frac{E}{2\left(r+\frac{R}{2}\right)}$，

	导体棒对挡条的压力为零，由平衡条件得

	$B_2I_{ab}d=mg$ 或 $B_2=\frac{mg(R+2r)}{Ed}$，

	联立解得 $B_2=0.50\UT$，方向垂直纸面向外。

	（2）棒下滑的过程中，由动量定理得

	$(mg-\bar{I}B_2d)t=mv$ 或 $mgt-B_2d\Delta q=mv$，

	又 $\Delta q=\bar{I}t=\frac{dhB_2}{2R}$，

	联立解得 $v=gt-\frac{hB_2^{2}d^{2}}{2Rm}$，

	由功能关系知，ab 棒产生的热量为

	$Q=\frac{1}{2}\left(mgh-\frac{1}{2}mv^{2}\right)$，

	代入解得 $Q=2.3\times10^{-3}\UJ$。
}

\item
%% number: 23
%% paperName: 2017年11月浙江选考第 23 题 10 分
%% typeId: 6
%% type: 计算题
%% chapter: 磁场
%% point: 带电粒子在磁场中的运动
%% method: 无
%% score: 10
%% degree: 350
%% duplicateId: 0
%% body:
如图所示，$x$ 轴上方存在垂直纸面向外的匀强磁场，坐标原点处有一正离子源，单位时间在 $xOy$ 平面内发射 $n_0$ 个速率为 $v$ 的离子，分布在 $y$ 轴两侧各为 $\theta$ 的范围内。在 $x$ 轴上放置长度为 $L$ 的离子收集板，其右端点距坐标原点的距离为 $2L$，当磁感应强度为 $B_0$ 时，沿 $y$ 轴正方向入射的离子，恰好打在收集板的右端点。整个装置处于真空中，不计重力，不考虑离子间的碰撞，忽略离子间的相互作用。
\begin{enumerate}
	\item
	求离子的比荷 $\frac{q}{m}$；
	\item
	若发射的离子被收集板全部收集，求 $\theta$ 的最大值；
	\item
	假设离子到达 $x$ 轴时沿 $x$ 轴均匀分布。当 $\theta=37^\circ$，磁感应强度在 $B_0\leq B\leq3B_0$ 的区间取不同值时，求单位时间内收集板收集到的离子数 $n$ 与磁感应强度 $B$ 之间的关系（不计离子在磁场中运动的时间）。
\end{enumerate}

\onepicture[w=7.0cm, align=right]{figs/23.png}
%% answer: 
\jdanswer{
\begin{enumerate}
	\item
	$\frac{v}{B_0L}$

	\item
	$\frac{\pi}{3}$

	\item
	$B_0\leq B\leq1.6B_0$ 时，$n_1=n_0$；$1.6B_0<B\leq2B_0$ 时，$n_2=n_0\left(5-\frac{5B}{2B_0}\right)$；$2B_0<B\leq3B_0$ 时，$n_3=0$。

\end{enumerate}
}
%% memo:
\memoanswer{
	（1）由洛伦兹力提供向心力得

	$qvB_0=m\frac{v^{2}}{R}$，

	又圆周运动的半径为

	$R=L$，

	联立解得 $\frac{q}{m}=\frac{v}{B_0L}$。

	（2）如图甲所示，以最大值 $\theta_m$ 入射时，由几何关系得

	$\Delta x=2R(1-\cos\theta_m)=L$ 或 $2R\cos\theta_m=L$，

	解得 $\theta_m=\frac{\pi}{3}$。

	\onepicture[w=5.5cm]{figs/23_an1.jpg}

	（3）当 $B>B_0$ 时，全部收集到离子时的最小半径为 $R_1$，如图乙所示，由几何关系得

	$2R_1\cos37^\circ=L$，解得 $B_1=\frac{mv}{qR_1}=1.6B_0$，

	当 $B_0\leq B\leq1.6B_0$ 时，有 $n_1=n_0$，

	当 $B>1.6B_0$ 时，恰好收集不到离子时的半径为 $R_2$，由几何关系得

	$R_2=0.5L$，解得 $B_2=2B_0$，

	当 $1.6B_0<B\leq2B_0$ 时，设 $R'=\frac{mv}{qB}$，有

	$n_2=\frac{2R'-L}{2R'(1-\cos37^\circ)}n_0=n_0\left(5-\frac{5B}{2B_0}\right)$，

	当 $2B_0<B\leq3B_0$ 时，有 $n_3=0$。

	\onepicture[w=5.5cm]{figs/23_an2.jpg}
}

\end{enumerate}

\end{document}
